% =================================================
% Ученический рабочий лист
% =================================================

\worksheettitle
  {Ломаная и многоугольник}
  {Звенья, вершины, стороны и диагонали}

\studentnameblock

\begin{checkpointbox}[Входная разминка]
  Отметьте три точки \(A\), \(B\), \(C\), не лежащие на одной прямой.
  Проведите отрезки \(AB\), \(BC\) и \(CA\). Какую фигуру вы получили?
  \answerblank[55mm]
\end{checkpointbox}

\section{Ломаная}

\openbrokenfigure

\begin{fillbox}[Дополните определение]
  Ломаная состоит из последовательно соединённых
  \answerblank[34mm]. Эти отрезки называют \answerblank[28mm], а их
  концы --- \answerblank[32mm] ломаной.
\end{fillbox}

\begin{practicebox}[Задание 1. Прочитайте рисунок]
  Для ломаной \(ABCDE\):
  \begin{enumerate}[label=\alph*),itemsep=0.5em]
    \item число вершин: \answerblank[18mm];
    \item число звеньев: \answerblank[18mm];
    \item перечислите все звенья: \answerblank[90mm].
  \end{enumerate}
\end{practicebox}

\section{Замкнутая ломаная}

\studentrecognizefigure

\begin{practicebox}[Задание 2. Распознайте фигуры]
  \begin{enumerate}[label=\alph*),itemsep=0.55em]
    \item Какие рисунки изображают незамкнутые ломаные?
      \answerblank[45mm]
    \item Какие рисунки изображают замкнутые ломаные?
      \answerblank[45mm]
    \item На каком рисунке проведена дополнительная диагональ?
      \answerblank[20mm]
  \end{enumerate}
\end{practicebox}

\section{Многоугольник}

\polygonpartsfigure

\begin{fillbox}[Главные слова]
  Звенья границы многоугольника называют \answerblank[32mm], а их концы ---
  \answerblank[32mm]. Число сторон многоугольника \answerblank[22mm]
  числу его вершин.
\end{fillbox}

\begin{practicebox}[Задание 3. Назовите фигуру]
  У изображённого многоугольника:
  \begin{enumerate}[label=\alph*),itemsep=0.5em]
    \item число вершин: \answerblank[18mm];
    \item число сторон: \answerblank[18mm];
    \item название по числу сторон: \answerblank[45mm];
    \item перечислите стороны: \answerblank[82mm].
  \end{enumerate}
\end{practicebox}

\polygonnamesfigure

\begin{practicebox}[Задание 4. Заполните таблицу]
  \begin{center}
    \renewcommand{\arraystretch}{1.45}
    \begin{tabularx}{\textwidth}{>{\centering\arraybackslash}p{0.2\textwidth}
      >{\centering\arraybackslash}X >{\centering\arraybackslash}X
      >{\centering\arraybackslash}X}
      \toprule
      Число сторон & 3 & 5 & 6 \\
      \midrule
      Название & \answerblank[30mm] & \answerblank[30mm] & \answerblank[30mm] \\
      \bottomrule
    \end{tabularx}
  \end{center}
\end{practicebox}

\section{Название многоугольника}

\correctnamingfigure

\Needspace{18\baselineskip}
\begin{thinkbox}[Задание 5. Какие названия верны?]
  Отметьте все правильные названия изображённого многоугольника:
  \[
    \square\ ABCDEF\qquad
    \square\ BCDEFA\qquad
    \square\ ACDEFB\qquad
    \square\ AFEDCB.
  \]
  Объясните, почему одно из названий неверно.
  \writinglines{2}
\end{thinkbox}

\section{Диагонали}

\diagonalfigure

\begin{fillbox}[Дополните]
  Диагональ соединяет две \answerblank[34mm] вершины многоугольника.
  Из вершины \(A\) на рисунке проведены диагонали
  \answerblank[65mm].
\end{fillbox}

\clearpage
\begin{practicebox}[Задание 6. Диагонали шестиугольника]
  Проведите все диагонали шестиугольника, выходящие из вершины \(A\).
  \hexagondiagonalsstudent
  Их количество: \answerblank[20mm].
\end{practicebox}

\section{Многоугольники внутри многоугольника}

\begin{practicebox}[Задание 7. Построение]
  Начертите пятиугольник \(ABCDE\). Проведите диагональ \(AC\). Назовите
  два образовавшихся многоугольника.
  \studentblankpentagon
  Ответ: \answerblank[55mm] и \answerblank[65mm].
\end{practicebox}

\clearpage
\begin{thinkbox}[Задание 8. Более сложная схема]
  В четырёхугольнике \(ABCD\) точка \(M\) --- середина стороны \(AD\).
  Проведите \(AC\) и \(BM\). Точку их пересечения обозначьте \(P\).
  \studentquadrilateralfigure
  Назовите все образовавшиеся части:
  \writinglines{2}
\end{thinkbox}

\begin{checkpointbox}[Задание 9. Верно или неверно]
  Поставьте знак «+» или «−».
  \begin{enumerate}[label=\alph*),itemsep=0.5em]
    \item У пятиугольника пять вершин. \answerblank[12mm]
    \item Любой отрезок между двумя вершинами является диагональю. \answerblank[12mm]
    \item Стороны \(AB\) и \(BA\) --- одна и та же сторона. \answerblank[12mm]
    \item В названии многоугольника вершины можно перечислять в любом порядке.
      \answerblank[12mm]
    \item Из одной вершины семиугольника выходят четыре диагонали.
      \answerblank[12mm]
  \end{enumerate}
\end{checkpointbox}

\begin{reflectionbox}[Итог]
  Чтобы правильно назвать многоугольник, нужно перечислять его вершины
  \answerblank[85mm].
\end{reflectionbox}
